\documentclass[10pt]{article}
\usepackage[a4paper,margin=24mm]{geometry}
\usepackage{amsmath,amssymb}
\usepackage[T1]{fontenc}
\usepackage{lmodern}
\usepackage{hyperref}
\title{A one-station counterexample to Conjecture 00000005508}
\author{ziangni-sys}
\date{4 October 2026}
\begin{document}
\maketitle
\section*{The clause being disproved}
The original conjecture defines a bottleneck as the first station to saturate
and asserts that the number of tied bottlenecks is the multiplicity of the
minimum residual capacity, and is at most the logarithm of the station count.
The Chinese version states the same bound on that multiplicity.
There is no restriction excluding a network with one station.
We disprove this bound, and therefore the conjunction, without needing to
assess its identification or switching clauses.

\section*{Network and proof}
Consider a network consisting of one station with service capacity
$\mu_0=1$. Jobs enter at rate $t\geq0$ and leave after service.
Its traffic coefficient is $a_0=1$, its offered traffic is $ta_0=t$,
and its residual capacity is
\[
r_0(t)=\mu_0-ta_0=1-t.
\]
Saturation means that offered traffic reaches or exceeds capacity:
$\mu_0\leq ta_0$. The station is unsaturated for every $0\leq t<1$
and saturates at $t=1$. It is therefore the first station to saturate.
The bottleneck set is the singleton $\{0\}$ and has cardinality one.
The residual-capacity vector has one entry, so its minimum has
multiplicity one at every load, including the saturation load.
We use precisely the multiplicity convention explicitly specified in the
conjecture; it does not assign multiplicity zero to a unique minimizer.

The station count is $n=1$. The bound would read
\[
1\leq\log1=0,
\]
which is false. This is independent of logarithm base: for every genuine
base $b>0$, $b\neq1$, $\log_b1=0$.
For comparison, a network of two identical independent stations receiving
rate $t$ each has residual vector $(1-t,1-t)$ and two tied bottlenecks at
$t=1$. It also violates the bound for natural logarithms
($\log2<1<2$), base two ($\log_22=1<2$), and base ten
($\log_{10}2<1<2$). The singleton witness additionally resolves every
possible legitimate base, without relying on a plurality convention.
The claim is an exact bound; interpreting it as an asymptotic order bound
or adding an additive constant would change the conjecture.

\section*{Formal verification and scope}
The Lean project uses Mathlib's real numbers and natural logarithm.
\texttt{Network n} records positive capacities and positive traffic
coefficients on the actual station type \texttt{Fin n}.
\texttt{Saturated} is the capacity inequality above.
\texttt{FirstBottleneck} requires a nonnegative load at which the station
saturates and the absence of saturation at every earlier nonnegative load,
for every station. The bottleneck set is the filtered finite set of stations
satisfying this predicate, and its cardinality is computed from that set.

\texttt{single\_first} proves first saturation at load one.
\texttt{single\_tied\_count} proves that the defined bottleneck set has
cardinality one. The minimum multiplicity is computed independently by
the minimum-station set; the two counts agree at saturation.
The final theorem
\texttt{conjecture\_5508\_false} refutes the universally quantified
logarithmic tie bound.
\texttt{single\_violates\_every\_base} verifies the base-change expression
$\log1/\log b$ as well.
The two-station comparison is supplementary; the fully formalized singleton
witness alone disproves the exact original claim.
No statistical assumption or queueing simulation is needed: this is a
deterministic traffic/capacity network.
\end{document}
